% !PRINTED 778
\weberpage{778}
% running head "II. Nachtrag." at the top of this page is not transcribed as body text.

Wir haben aber hier nicht nöthig, von den speciellen Ausdrücken dieser
Charaktere Gebrauch zu machen. Dahingegen stützen wir uns auf folgenden
Lehrsatz, der für alle \gesperrt{Abel}'schen Gruppen gilt:

1. \gesperrt{Ist} $A$ \gesperrt{ein Element} $f^{\text{ten}}$ \gesperrt{Grades
einer} \gesperrt{Abel}'\gesperrt{schen Gruppe} $\mu^{\text{ten}}$
\gesperrt{Grades, und ist} $\mu = ef$, \gesperrt{so sind alle} $\chi_i(A)$
$f^{\text{te}}$ \gesperrt{Einheitswurzeln, und darunter kommt jede}
$f^{\text{te}}$ \gesperrt{Einheitswurzel genau} $e$ \gesperrt{mal vor.}

Der Satz ist implicite in dem Satze 7., \S.\,12 dieses Bandes enthalten. Denn
wenden wir diesen Satz auf die Gruppe
\[
T = 1, A, A^2, \ldots, A^{f-1}
\]
an, deren Index in Bezug auf $\mathfrak{N}$ gleich $e$ ist, so folgt, dass es
eine Gruppe $\Xi$ von $e$ Charakteren $\xi$ giebt, die den
Bedingungen
\[
\xi(A) = 1
\]
genügen.

Zerlegt man also die Gruppe $X$ der Charaktere $\chi$ in die Nebengruppen
\[
\Xi_1, \Xi_2, \ldots, \Xi_f,
\]
so sind $\chi_1(A)$ und $\chi_2(A)$ einander gleich oder von einander
verschieden, je nachdem $\chi_1$ und $\chi_2$ in derselben oder in
verschiedenen dieser Nebengruppen vorkommen. Da ausserdem wegen
$\chi(A)^f = \chi(A^f) = 1$ alle $\chi(A)$ Einheitswurzeln vom Grade $f$
sind, und es nur $f$ verschiedene solche giebt, so ist damit unser Theorem
bewiesen.

Ist $n$ in der Classe $A$ enthalten, so ist der Grad $f$ von $A$ der
Exponent, zu dem $n$ nach dem Modul $m$ \gesperrt{gehört}, d.\,h. der kleinste positive
Exponent, für den
\weq{5}{n^f \equiv 1 \pmod{m}}
ist. Wenn also $n$ zum Exponenten $f$ gehört, so sind alle $\chi_i(n)$
$f^{\text{te}}$ Einheitswurzeln, und jede $f^{\text{te}}$ Einheitswurzel
kommt darunter $e$ mal vor.

Bedeutet $a$ irgend einen der $\varphi(m)$ Reste der Zahlen $n$, so ist in der
Form
\weq{6}{\mu_x = mx + a - m}
jede Zahl aus der durch $a$ repräsentirten Zahlclasse (2) enthalten;
verstehen wir unter $a$ den kleinsten positiven Rest, so erhalten wir die
Gesammtheit der Zahlen dieser Classe, wenn wir $x$ alle positiven
ganzzahligen Werthe durchlaufen lassen.
